\scanpage{313}
\mspage{12}
% archive sheet number: 314

Alors la connaissance de $C_{\cdot\deleted{\unsure{\ast}}}=\mathbb C_{k\cdot}(X_{\ast})
=\mathbb C_{\cdot}\otimes_{\mathbb Z}k$ est encore suffisante pour
la construction de $F(X_{\ast})$, car on aura
maintenant
\[
\begin{aligned}
F(X_{\ast})&\simeq\operatorname{Hom}_{k\cdot}(C_{\cdot k},K_{\cdot})\\
&\simeq Z^0\bigl(\underbrace{\mathcal{H}\mathrm{om}^{\cdot}_{k}{}_{\cdot}(C_{\cdot k},K_{\cdot})}_{
\begin{gathered}\mathcal{H}\mathrm{om}^{i}_{k}{}_{j}(C_{\cdot k}^{\unsure{\ell|i}},K_{\cdot})\\
=\operatorname{Hom}_{k}(C_{ik},K_j))\end{gathered}}\bigr)
\end{aligned}
\]

($\mathrm{Hom}_{k}(A,M)$, pour $A\in\mathit{Ab}_{k}$ et $M\in B$, étant
défini par
\[
\operatorname{Hom}_{B}(N,\mathrm{Hom}_{k}(A,M))\simeq
\operatorname{Hom}_{k}(A,\operatorname{Hom}_{B}(N,M))).
\]

2°) Si on regarde toutes les façons possibles
\deleted{d'envoyer} \inserted{d'envoyer}
\[
(\mathit{Ss})^{\circ}\xrightarrow{F}B
\]
dans $B$ stable par $\varprojlim$, avec $F$ y
commutant, il y en a \deleted{bien sûr} une
``universelle'', dont toutes les autres sont
``spécialisation'', obtenue en prenant
\[
B_0=(\mathit{Ab}_{\cdot})^{\circ},\quad\text{et}\quad
F_0:(\mathit{Ss})^{\circ}\longrightarrow B_0=(\mathit{Ab}_{\cdot})^{\circ}
\]
défini par
\[
F_0(X_{\ast})=\mathbb C_{\cdot}(X_{\ast})
\]
(qui dépend de $X_{\ast}$ de façon contravariante
quand $\mathbb C_{\cdot}(X_{\ast})$ \unsure{s'entend alors|s'interprète|s'envisage} comme
objet de $(\mathit{Ab}_{\cdot})^{\circ}$). De façon précise ;

\uline{Prop. 3} Considérons
% margin note written sideways
\margin{pourrait passer pour l'essentiel avant le \unsure{topos|topo} de \uline{Dold-Puppe}}
\[
\mathit{Ss}^{\circ}\xrightarrow{F_0}B_0=(\mathit{Ab}_{\cdot})^{\circ}
\simeq(\mathit{Ab}^{\circ})^{\cdot}
\]

défini par $F_0(X_{\ast})=\mathbb C_{\cdot}(X_{\ast})$. Alors
$B_0$ est une catégorie \inserted{add} avec $\varprojlim$, et
pour tte autr. telle catégorie $B$,
le foncteur
\[
\underline{\operatorname{Hom}}_{\substack{\mathrm{Comm.}\\\mathrm{aux}\ \varprojlim}}
(\deleted{\unsure{\mathit{Ss}}}B_0,B)\longrightarrow
\underline{\operatorname{Hom}}_{\substack{\mathrm{comm}\\\varprojlim}}(\mathit{Ss}^{\circ},B)
\]
donné par $\varphi\longmapsto\varphi\circ F_0$
\margin{\unsure{21} ($B_{\cdot}$ par \unsure{pr.|prop.} 2)}
est une équivalence de catégories.

\deleted{\unsure{Compte tenu de}
Compte tenu de prop 2, cela signifie
de f.\ \unsure{aussi} que la}
Grâce au \unsure{théor.} de Dold-Puppe, on peut
\deleted{Cela signifie que $F_0$ établit l'\unsure{iso}
Grâce aux : l'équivalence \unsure{coeff.ts} de degré
\unsure{purement} \unsure{=}. Ou bien dire : indiquer la condition
de comm. aux $\varprojlim$)
$\deleted{\unsure{\operatorname{Hom}}}\ \underline{\operatorname{Hom}}_{!}(B_0,B)$}

remplacer dans la déf. de
$F_0,B_0$, $(\mathit{Ab}_{\cdot})^{\circ}$ par $(\mathit{Ab}_{\ast})^{\circ}=\widehat\Delta_{\mathrm{ab}}^{\circ}$,
$\mathbb C_{\cdot}(X_{\ast})$ par $\mathbb C_{\ast}(X_{\ast})$\inserted{$=\mathbb Z(X_{\ast})$}, et l'équivalence ci-dessus
devient
\[
\underline{\operatorname{Hom}}_{\deleted{\unsure{\$}}}^{!}(\widehat\Delta_{\mathrm{ab}}^{\circ},B)
\overset{\overset{?}{\simeq}}{\longrightarrow}
\underline{\operatorname{Hom}}_{\unsure{\$}}^{!}(\widehat\Delta^{\circ},B)
\]

où l'\unsure{explosion|exposant|expression} $!$ signifie qu'on se borne
aux foncteurs qui commutent aux $\varprojlim$.
Or sous cette forme, l'assertion est
valable quelle que soit la cat. (ens. \unsure{petite|pointée})
$\Delta$, \deleted{la démonstration se faisant comme ci-dessus.}
\inserted{On a un foncteur ``abélianisation''
$X\longrightarrow\mathbb Z(X_{\cdot})\qquad\widehat\Delta\longrightarrow\widehat\Delta_{\mathrm{ab}}$
commutant aux $\varinjlim$ de \deleted{$\widehat\Delta^{\circ}\longrightarrow B=(\unsure{\widehat{\mathit{Ab}}|\widehat{\Delta}_{\mathrm{ab}}})^{\circ}$}}
